% Part: proof-theory
% Chapter: proof-search
% Section: search-algorithm

\documentclass[../../../include/open-logic-section]{subfiles}

\begin{document}

\olfileid[pt]{pt}{ps}{sal}

\olsection{Algoritmo de busca para \Log{G3c}}

Descrevemos um algoritmo de busca para \Log{G3c}. Ele não é o único
possível, nem sequer o mais eficiente. Mas é correto: se um sequente
$\Gamma \Sequent \Delta$ tiver !!a{proof}, o algoritmo acabará parando
e encontrará !!a{proof}. Ele também é infalível: se parar sem ter
encontrado !!a{proof}, ou nunca parar, então $\Gamma \Sequent \Delta$
não era válido desde o início.

Uma característica necessária e essencial do algoritmo é garantir que
toda !!{formula} em $\Gamma \Sequent \Delta$, ou em um sequente
produzido durante a busca de provas, seja ``reduzida'', isto é, seja
tratada como a fórmula principal de uma regra de inferência aplicada de
trás para a frente, e que isso aconteça tantas vezes quanto necessário.
Chamamos essa característica de \emph{equidade}

O algoritmo funciona em estágios. A saída de cada estágio é uma árvore
de sequentes, e a saída do estágio seguinte é produzida reduzindo
!!a{formula} em cada sequente superior que ainda não seja um axioma.

No estágio $0$, começamos com $\Gamma \Sequent \Delta$. Para garantir a
equidade, atribuímos números como índices às !!{formula}s de $\Gamma$ e
$\Delta$, de modo que !!{formula}s diferentes recebam números
diferentes. Por exemplo, imagine contar as !!{formula}s da esquerda
para a direita. Escrevemos os números como sobrescritos. Por exemplo,
se estivermos buscando !!a{proof} de $!A, !B \Sequent !C$, a saída do
estágio~$0$ será $!A^1, !B^2 \Sequent !C^3$. O maior~$n$ que ocorre
como sobrescrito em um sequente indexado desse tipo é seu \emph{índice
máximo} ($3$~no exemplo).

Se o sequente contiver quantificadores, também precisaremos saber quais
termos usar ao reduzir $\lexists[x][!A(x)]$ à direita ou
$\lforall[x][!A(x)]$ à esquerda. Precisaremos apenas de termos
construídos a partir das !!{constant}s $\Obj c_0$, $\Obj c_1$, $\Obj
c_2$, \dots, usando quaisquer símbolos de função que ocorram
em~$\Gamma$ e $\Delta$. Supomos que o conjunto de todos esses termos
seja fornecido em alguma enumeração fixa, para que possamos falar, por
exemplo, da ``primeira !!{constant} que não ocorre em $\Gamma \Sequent
\Delta$''. Cada sequente indexado da árvore produzida pelo algoritmo tem
um conjunto associado de !!{constant}s. O conjunto de !!{constant}s
atribuído ao sequente no estágio~$0$ é o conjunto de todas as
!!{constant}s que ocorrem nele (se houver alguma), e $\{\Obj c_0\}$ se
ele não contiver !!{constant}s.

Suponha que o algoritmo tenha gerado, no estágio~$n$, uma árvore de
sequentes indexados com conjuntos associados de !!{constant}s. No
estágio~$n+1$, fazemos o seguinte para cada sequente superior.

Se houver !!a{formula}~$!A$ que ocorra tanto à esquerda quanto à direita
do sequente, ou se o lado esquerdo contiver~$\lfalse$, não fazemos nada:
tais sequentes têm !!{proof}s em~\Log{G3c}, e não resta nada a fazer
para esse sequente superior específico.

Se não houver nenhuma !!{formula} não atômica no sequente, aborte o
algoritmo. Não há prova de $\Gamma \Sequent \Delta$.

Caso contrário, escolha a !!{formula} não atômica~$!A$ com o menor
índice~$i$. Então, o sequente superior em questão é $!A^i, \Pi \Sequent
\Lambda$ ou $\Pi \Sequent \Lambda, !A^i$. Seja $k$ o índice máximo do
sequente e $C$ o conjunto de !!{constant}s atribuído a ele.

``Reduzimos'' $!A^i$, isto é, geramos novos sequentes superiores de
acordo com a regra de inferência determinada pelo !!{main operator}
de~$!A$ e conforme $!A^i$ apareça à esquerda ou à direita.

\begin{enumerate}
  \item $!A \ident !B \land !C$ e ocorre à esquerda: acrescente um novo
  sequente superior acima de $(!B \land !C)^i, \Pi \Sequent \Lambda$:
  \[
  \Axiom$!B^k, !C^{k+1}, \Pi \fCenter \Lambda$
  \RightLabel{\LeftR{\land}}
  \UnaryInf$(!B \land !C)^i, \Pi \fCenter \Lambda$
  \DisplayProof
  \]
  O conjunto de !!{constant}s atribuído ao novo sequente é~$C$.
  \item $!A \ident !B \land !C$ e ocorre à direita: acrescente dois
  novos sequentes superiores acima de $\Pi \Sequent \Lambda, (!B \land
  !C)^i$:
  \[
  \Axiom$\Pi \fCenter \Lambda, !B^k$
  \Axiom$\Pi \fCenter \Lambda, !C^k$
  \RightLabel{\RightR{\land}}
  \BinaryInf$\Pi \fCenter \Lambda, (!B \land !C)^i$
  \DisplayProof
  \]
  O conjunto de !!{constant}s atribuído ao novo sequente é~$C$.

  \item Os casos em que $!A \ident \lnot !B$, $!A \ident !B \lor !C$
  e $!A \ident !B \lif !C$ são semelhantes e ficam como exercícios.

  \item $!A \ident \lexists[x][!B(x)]$ e ocorre à esquerda: acrescente
  um novo sequente superior acima de $\lexists[x][!B(x)]^i, \Pi
  \Sequent \Lambda$:
  \[
  \Axiom$!B(c)^k, \Pi \fCenter \Lambda$
  \RightLabel{\LeftR{\lexists}}
  \UnaryInf$\lexists[x][!B(x)]^i, \Pi \fCenter \Lambda$
  \DisplayProof
  \]
  em que $c$ é a primeira !!{constant} que não pertence a~$C$. O
  conjunto de !!{constant}s atribuído ao novo sequente é $C \cup
  \{c\}$.
  \item $!A \ident \lexists[x][!B(x)]$ e ocorre à direita: acrescente um
  novo sequente superior acima de $\Pi \Sequent \Lambda,
  \lexists[x][!B(x)]$:
  \[
  \Axiom$\Pi \fCenter \Lambda, \lexists[x][!B(x)]^k, !B(t)^{k+1}$
  \RightLabel{\LeftR{\lexists}}
  \UnaryInf$\Pi \fCenter \Lambda, \lexists[x][!B(x)]^i$
  \DisplayProof
  \]
  em que $t$ é o primeiro termo que contém apenas !!{constant}s de~$C$
  e que ainda não foi usado na redução de $\lexists[x][!B(x)]$ à
  direita abaixo desse sequente (ou $\Obj c_0$ se todos esses termos já
  tiverem sido usados). O conjunto de !!{constant}s atribuído ao novo
  sequente é $C$.
  \item Os casos em que $!A \ident \lforall[x][!B(x)]$ são semelhantes
  e ficam como exercícios.
\end{enumerate}

Se, no estágio $n+1$, não tivermos acrescentado um novo sequente a
nenhum sequente superior, o algoritmo terminará com êxito. Nesse caso,
teremos encontrado !!a{proof} de~$\Gamma \Sequent \Delta$, obtida da
árvore no estágio $n+1$ apagando todos os índices. Todos os sequentes
superiores têm a forma $!A, \Pi \Sequent \Lambda, !A$ ou a forma
$\lfalse, \Pi \Sequent \Lambda$. Todas as inferências são inferências
corretas de~\Log{G3c}. Isso é evidente para as inferências
proposicionais. Observe que, em cada passo, o conjunto~$C$ de
!!{constant}s atribuído a um sequente contém todas as !!{constant}s que
ocorrem no sequente. Portanto, na redução que usa \LeftR{\lexists} e
\RightR{\lforall}, a condição de variável própria é satisfeita, pois a
variável própria~$c$ era !!a{constant} que não pertencia ao conjunto~$C$
atribuído à conclusão; $c$ também não ocorria na conclusão. Temos:

\begin{prop}
O algoritmo de busca de provas para \Log{G3c} é correto: se terminar
com êxito, o sequente inicial~$\Gamma \Sequent \Delta$ terá !!a{proof}.
\end{prop}





\end{document}
